\section{Acquired strata with changing rank}\label{sec:singular}
The preceding quantum class stays in the positive definite interior. We now allow a physical rank change while retaining its actual probability. The tool is the curvature alternative, which does not require the same-report map to contract on a boundary stratum.

\begin{proposition}[A weighted acquired-stratum profile]\label{prop:strata}
Suppose the affine certificate and (C) hold on a compact convex domain $K$, with $b_t$ satisfying its product bound. Suppose each actual row law has the decomposition
\[
 \mu_{t,u}=\sum_{j=1}^J w_j\mu_{t,u}^{(j)},\qquad
 w_j>0,\quad\sum_jw_j=1,
\]
where, uniformly in $t,u$, the probability $\mu_{t,u}^{(j)}$ obeys a small-ball upper bound of exponent $d_j>0$. Its support lies in a common compact set $E_{t,j}\subset K$ with at-most-$m$ Euclidean covers of radius $C_jm^{-1/d_j}$. The component mark is visible in the report, or there is a common Borel assignment using the corresponding union codebook. Put
\begin{equation}\label{eq:profile}
 \Xi_w(m)=\sum_{j=1}^Jw_jm^{-2/d_j}.
\end{equation}
Then the checkpoint excess is comparable to $\Xi_w(M)$, and under the common-support autonomous interface its excess is comparable to $\Xi_w(M-n)$ for every $M\ge n+1$. Constants may depend on $J$, the displayed geometric and stability constants, but need not discard or normalize any $w_j$. If atoms are present their exact finite supports may instead be assigned a declared finite number of labels and zero residual profile thereafter.
\end{proposition}
\begin{proof}
For any global $m$-center set, its distortion on component $j$ is at least $q_m(\mu^{(j)})$. Summing gives the lower $c\Xi_w(m)$, with no assumption that different components use disjoint centers. The actual full-history mixture has the same component weights and small-ball bounds, since these are uniform in its incoming prediction.

For $m\ge2J$, allocate $\lfloor m/J\rfloor$ centers to each component; their union uses at most $m$ centers and has distortion at most $C\Xi_w(m)$ under every row. For $m<2J$, a single legal response has bounded error and $\Xi_w(m)\ge(2J)^{-a_*}$, where $a_*=\max_j2/d_j$, so enlargement of a fixed constant gives the same upper. The representative after pooling lies in $K$, even if it lies on none of the acquired strata. Equation~\eqref{eq:pool-cover} applies to the union codebook and gives
$r_t\le b_t r_{t-1}+C\Xi_w(m_t)$.

The clock lower and row lower yield $c\Xi_w(M-n)$. In the allocation~\eqref{eq:allocation} choose $z=\rho^{1/(a_*+1)}$. Each geometric series $\sum_{k\ge0}\rho^kz^{-2k/d_j}$ is bounded. The recursion thus gives $C\sum_jw_j(M-n)^{-2/d_j}$, including the minimal-budget case as before. Finite atoms can be covered exactly; the additional labels must be counted rather than deleted from the budget. This proves the stated extension.
\end{proof}

Here is a raw instrument with two genuinely different acquired ranks. Let $0<\lambda<1$ and first apply
$q=T(\rho)=\lambda\rho+(1-\lambda)I_D/D$ to an arbitrary density operator, possibly singular. With probability $\omega\in(0,1)$ perform the weak effect instrument of Section~\ref{sec:quantum} \emph{after} this channel, with unnormalized output $K_HqK_H\,\mu(dH)$. With probability $1-\omega$ perform the rank-one instrument
\begin{equation}\label{eq:pure-instrument}
 q\longmapsto D\operatorname{tr}(P_vq)P_v\,\sigma(dv),
 \qquad P_v=vv^*,
\end{equation}
where $\sigma$ is the unitary-invariant probability on complex projective space. The branch mark is part of the report. The instrument is completely positive: its density is the map $q\mapsto D P_vqP_v$. Since $\int D P_v\,d\sigma=I_D$, both branches are normalized separately. The common preparation channel is charged within each instrument call; it is not a free regeneration or a discard-and-repeat protocol.

\begin{theorem}[Repeated full-rank/pure-rank acquisition]\label{thm:singular}
Equip this raw instrument with the same separating audit as in~\eqref{eq:audit-normalization}. If
\begin{equation}\label{eq:singular-curvature}
 b_*:=\omega\lambda^2\frac{(1+2\eta)^2}{1-\eta}<1,
\end{equation}
then, for $n\ge1$, the actual predictive law has its full-rank and rank-one components with masses exactly $\omega$ and $1-\omega$. With $d_f=D^2-1$ and $d_p=2(D-1)$,
\begin{align}
 \calR_{\rm cp}(n,M)&=B_n+\Theta\bigl(\omega M^{-2/d_f}+(1-\omega)M^{-2/d_p}\bigr),\label{eq:singular-cp}\\
 \calR_{\rm aut}(n,M)&=B_n+\Theta\bigl(\omega(M-n)^{-2/d_f}+(1-\omega)(M-n)^{-2/d_p}\bigr)
 \label{eq:singular-aut}
\end{align}
for feasible $M\ge n+1$. Constants have the dependencies in Proposition~\ref{prop:strata}, including the distance of $b_*$ from one. This is a total-risk statement with the actual baseline $B_n$, not the baseline of the nonsingular experiment.
\end{theorem}
\begin{proof}
Every $q$ has eigenvalues at least $(1-\lambda)/D$, uniformly over the entire closed state body. In the weak branch Lemma~\ref{lem:congruence} therefore supplies a bounded density of the full-rank output in $d_f$ affine coordinates. Its support lies in a compact full-rank subset with a global $d_f$-dimensional cover. The pure branch has density $D\operatorname{tr}(P_vq)\le D$ relative to $\sigma$. For a fixed pure projector $P_w$,
\[
 \|P_v-P_w\|_{\rm HS}^2=2(1-|v^*w|^2),\qquad
 \sigma\{\|P_v-P_w\|_{\rm HS}<h\}=(h^2/2)^{D-1}
\]
for $0\le h\le\sqrt2$. To verify the cap formula, normalize independent standard complex Gaussian coordinates. Their squared moduli are independent exponential variables; integrating their sum shows that the first normalized squared modulus has density $(D-1)(1-u)^{D-2}$ on $[0,1]$. For a ball centered off the projective orbit, choose any orbit point in it and enlarge the radius by two. This gives the required small-ball bound for arbitrary centers.

For the global pure-state cover, select an index with $|v_j|\ge D^{-1/2}$, choose its phase positive, and use the other $D-1$ complex coordinates as a finite chart. On a slightly enlarged such chart the projector map is Lipschitz and has $2(D-1)$ real coordinates. A finite collection of chart grids yields $Ch^{-d_p}$ codepoints. No uniform simple spectrum or extraneous ambient exponent is used. The actual row measure is the indicated mixture with the original weights, because the branch coin is normalized before conditioning on its report. Full-history mixing over old predictions preserves those same weights and bounds.

For stability, use traceless tangent directions $v$. In the weak branch, with $u_H=K_HqK_H$ and $\ell_H=\operatorname{tr}(A_Hq)$,
\[
 \|D u_H[v]\|_{\rm HS}\le\lambda(1+\eta)\|v\|_{\rm HS},\qquad
 |D\ell_H[v]|\le\lambda\eta\|v\|_{\rm HS},\qquad
 \ell_H\ge1-\eta.
\]
A normalized density operator has Hilbert--Schmidt norm at most one. Equation~\eqref{eq:perspective-hessian} therefore bounds half the weak-branch Hessian by $\lambda^2(1+2\eta)^2/(1-\eta)$ times $\|v\|_{\rm HS}^2$. Centering by $I_D/D$ only subtracts an affine function from the potential and leaves its Hessian unchanged. In the pure branch the squared centered norm is the constant $1-1/D$, so its integrated potential is affine (in fact constant); its Hessian vanishes. The mixture has the bound~\eqref{eq:singular-curvature}. Derivatives extend to a neighborhood of the closed density body because $q$ and the effects retain strict positive lower bounds.

The report density in either component is affine and uniformly positive on its component support; their mixture, with the visible mark and fixed $\omega$, has common null sets at every phase. Total-variation continuity follows directly from the affine formulas. Proposition~\ref{prop:strata} now applies on the full convex density body. The true acquired rank changes between $D$ and one, while the coarse observer's conditional means may have any rank; all of them are included in the curvature certificate. The physical baseline $B_n=\E\|Z_n-X_n\|^2$ is kept throughout.
\end{proof}
